% Partie : théorie de la démonstration
% Chapitre : élimination des coupures
% Section : séquent médian

\documentclass[../../../include/open-logic-section]{subfiles}

\begin{document}

\olfileid[fr]{pt}{cut}{mid}

\olsection{Théorème du séquent médian et théorème de Herbrand}

Le théorème du séquent médian est une conséquence importante de
l'élimination des coupures. Il affirme que, s'il existe une
!!{proof}~$\pi$ de~$\Gamma \Sequent \Delta$ dont toutes les formules
du séquent final sont prénexes, il existe une !!{proof}~$\pi'$
contenant un séquent $S \ident \Pi \Sequent \Lambda$, appelé
\emph{séquent médian}, tel que les inférences situées au-dessus de~$S$
soient propositionnelles et celles situées au-dessous soient
des inférences de quantification.

Une !!{formula}~$!A$ est \emph{prénexe}, ou \emph{en forme prénexe},
si elle est de la forme
\[\mathsf{Q}_1x_1\dots\mathsf{Q}_nx_n\,!B(x_1,\dots,x_n)\]
où chaque $\mathsf{Q}_i$ est $\lforall$ ou~$\lexists$,
et où la matrice est sans quantificateurs.

Le théorème du séquent médian a pour conséquence importante le
théorème de Herbrand. Son cas général est plus délicat à formuler.
Schématiquement, à tout !!{sentence} démontrable~$!A$ de logique
du premier ordre correspond une tautologie propositionnelle~$!A'$
permettant de démontrer $!A$ par les instanciations et inférences
de quantification appropriées. Voici la version pour les
!!{sentence}s de la forme~$\lexists[x][!B(x)]$, avec $!B$
sans quantificateurs.

\begin{thm}[Théorème de Herbrand]
Dans le calcul sans égalité logique, supposons que $!B(x)$ soit sans
quantificateurs et que $\Proves \lexists[x][!B(x)]$. Il existe alors une
liste non vide de termes clos~$t_1$, \dots, $t_n$ telle que
$\Proves !B(t_1) \lor \dots \lor !B(t_n)$.
\end{thm}

La !!{formula}~$!B(t_1) \lor \dots \lor !B(t_n)$ est appelée
\emph{disjonction de Herbrand} de~$\lexists[x][!B(x)]$.
Puisque $!B$ est sans quantificateurs, cette disjonction l'est
également. Si elle est !!{provable}, l'élimination des coupures
et la propriété des sous-formules fournissent une preuve sans
coupure ne comportant que des inférences propositionnelles.
Elle est donc une tautologie.

La difficulté du théorème de Herbrand est d'établir l'existence
d'une disjonction de Herbrand pour toute formule démontrable
$\lexists[x][!B(x)]$. La réciproque est facile : si
$\lexists[x][!B(x)]$ possède une disjonction de Herbrand
démontrable, cette formule existentielle est démontrable.
La disjonction entraîne en effet $\lexists[x][!B(x)]$.
Voici une preuve dans~$\Log{G3c}$ lorsque $n=2$ ;
les deux séquents du haut sont dérivables par identité et
affaiblissement :
\[
\Axiom$!B(t_1) \fCenter \lexists[x][!B(x)], !B(t_1), !B(t_2)$
\Axiom$!B(t_2) \fCenter \lexists[x][!B(x)], !B(t_1), !B(t_2)$
\RightLabel{\LeftR{\lor}}
\BinaryInf$!B(t_1) \lor !B(t_2) \fCenter \lexists[x][!B(x)], !B(t_1), !B(t_2)$
\RightLabel{\RightR{\lexists}}
\UnaryInf$!B(t_1) \lor !B(t_2) \fCenter \lexists[x][!B(x)], !B(t_2)$
\RightLabel{\RightR{\lexists}}
\UnaryInf$!B(t_1) \lor !B(t_2) \fCenter \lexists[x][!B(x)]$
\DisplayProof
\]

Pour extraire une disjonction de Herbrand d'une !!{proof}~$\pi$
de~$\Sequent \lexists[x][!B(x)]$, considérons d'abord le cas
où toutes les inférences $\RightR{\lexists}$ de la preuve sans
coupure~$\pi$ sont regroupées à la fin :
\[
\AxiomC{}
\RightLabel{$\pi_1$}
\Deduce$ \fCenter \lexists[x][!B(x)], !B(t_1), \dots, !B(t_n)$
\RightLabel{\RightR{\lexists}}
\UnaryInf$ \fCenter \lexists[x][!B(x)], !B(t_2), \dots, !B(t_n)$
\RightLabel{$(n-1) \times \RightR{\lexists}$}
\Deduce$ \fCenter \lexists[x][!B(x)]$
\DisplayProof
\]
Si ce sont toutes les inférences $\RightR{\lexists}$ de~$\pi$
qui portent sur $\lexists[x][!B(x)]$, la sous-preuve~$\pi_1$
ne contient aucune autre inférence quantifiée : $!B(x)$ n'a pas
de quantificateurs, et le séquent final n'a aucune formule à gauche.
Ainsi, $ \fCenter \lexists[x][!B(x)], !B(t_1),
\dots, !B(t_n)$ est le séquent médian de~$\pi$.

Au-dessus de ce séquent, il n'y a pas d'inférence quantifiée.
Les occurrences de $\lexists[x][!B(x)]$ n'y sont donc ni actives
ni principales. Supprimons les occurrences de cette formule
dans la sous-preuve~$\pi_1$ : nous obtenons une preuve de
$\fCenter !B(t_1), \dots, !B(t_n)$, puis, par $n-1$ applications
de $\RightR{\lor}$, une preuve de la disjonction
$\Sequent !B(t_1) \lor \dots \lor !B(t_n)$.

Cependant, même dans une preuve sans coupure de
$\Sequent \lexists[x][!B(x)]$, les inférences
\RightR{\lexists} ne sont pas nécessairement regroupées à la fin.
Entre elles peuvent s'intercaler des inférences dont la formule
principale est une instance $!B(t_i)$.
D'où la nécessité du théorème du séquent médian.

\begin{thm}[Théorème du séquent médian]\ollabel{thm:midsequent}
S'il existe une !!{proof}~$\pi$ dans \Log{G3c} de
$\Gamma \Sequent \Delta$, où toutes les formules de $\Gamma$
et de $\Delta$ sont prénexes, il existe une !!{proof}~$\pi'$
contenant un séquent $\Pi \Sequent \Lambda$ tel que les inférences
au-dessous soient quantifiées et les inférences au-dessus
propositionnelles.
\end{thm}

\begin{proof}
Définissons l'\emph{ordre} d'une inférence quantifiée de~$\pi$
comme le nombre d'inférences propositionnelles situées au-dessous.
L'ordre~$o(\pi)$ de $\pi$ est la somme des ordres de ses inférences
quantifiées. Nous raisonnons par récurrence sur cet ordre.

Si $o(\pi) = 0$, aucune inférence quantifiée n'a d'inférence
propositionnelle au-dessous. Si la preuve ne contient aucune
inférence quantifiée, son séquent final est un séquent médian.
Sinon, comme chaque inférence quantifiée possède une seule
prémisse, ces inférences forment une chaîne finale. La prémisse
de la plus haute d'entre elles est le séquent médian recherché.

Si $o(\pi) > 0$, choisissons une inférence quantifiée d'ordre $>0$
la plus basse possible. Une inférence propositionnelle se trouve
au-dessous. La conclusion de l'inférence choisie ne peut être
la prémisse d'une autre inférence quantifiée : cette inférence
plus basse aurait elle aussi une inférence propositionnelle
au-dessous, contrairement au choix effectué.
Nous distinguons les cas selon les deux règles successives.

Puisque le séquent final ne contient que des formules prénexes,
la formule principale de l'inférence quantifiée ne peut être
active dans l'inférence propositionnelle suivante. Sinon,
la formule principale de cette dernière contiendrait un
quantificateur dans la portée d'un connecteur propositionnel.
La propriété des sous-!!{formula}s, prise avec les instances
de termes qu'autorise ce calcul, imposerait alors une telle
configuration dans une !!{formula} du séquent final.
La formule principale quantifiée appartient donc au contexte
de l'inférence propositionnelle inférieure.

Voici deux exemples.

Supposons d'abord que $\RightR{\lforall}$ aboutisse à la prémisse
droite de $\LeftR{\lif}$. La sous-!!{proof} pertinente,
notée~$\pi_1$, se termine ainsi :
\[
\AxiomC{}
\RightLabel{$\pi_2$}
\Deduce$\Gamma' \fCenter \Delta', \lforall[x][!A(x)], !B$
\AxiomC{}
\RightLabel{$\pi_3$}
\Deduce$!C, \Gamma' \fCenter \Delta', !A(c)$
\RightLabel{\RightR{\lforall}}
\UnaryInf$!C, \Gamma' \fCenter \Delta', \lforall[x][!A(x)]$
\RightLabel{\LeftR{\lif}}
\BinaryInf$!B \lif !C, \Gamma' \fCenter \Delta', \lforall[x][!A(x)]$
\DisplayProof
\]
Nous pouvons supposer toute la preuve~$\pi$ régulière.
La constante propre~$c$ n'apparaît donc pas en dehors de~$\pi_3$,
en particulier ni dans $!B \lif !C$ ni dans~$\pi_2$.
Considérons la preuve~$\pi_1'$ :
\[
\AxiomC{}
\RightLabel{$\pi_2'$}
\Deduce$\Gamma' \fCenter \Delta', \lforall[x][!A(x)], !A(c), !B$
\AxiomC{}
\RightLabel{$\pi_3[\Sequent \lforall[x][!A(x)]]$}
\Deduce$!C, \Gamma' \fCenter \Delta', \lforall[x][!A(x)], !A(c)$
\RightLabel{\LeftR{\lif}}
\BinaryInf$!B \lif !C, \Gamma' \fCenter \Delta', \lforall[x][!A(x)], !A(c)$
\RightLabel{\RightR{\lforall}}
\UnaryInf$!B \lif !C, \Gamma' \fCenter \Delta', \lforall[x][!A(x)], \lforall[x][!A(x)]$
\DisplayProof
\]
La preuve~$\pi_2'$ s'obtient à partir de~$\pi_2$ par affaiblissement
à droite avec~$!A(c)$. La prémisse issue de~$\pi_3$ est également
affaiblie à droite avec la formule universelle, comme l'indique
son annotation. Ces affaiblissements n'ajoutent aucune inférence
et ne modifient donc pas l'ordre. La régularisation préalable
permet d'éviter les conflits avec les constantes propres des
sous-preuves. La condition de constante propre de
\RightR{\lforall} reste satisfaite : $c$ n'apparaît pas
dans~$!B \lif !C$ ni dans le contexte final.

L'admissibilité de la contraction fournit alors une
!!{proof}~$\pi_1''$ de
$!B \lif !C, \Gamma' \fCenter \Delta', \lforall[x][!A(x)]$.
La preuve de l'admissibilité de $\RightR{\Contraction}$ pour
$\lforall[x][!A(x)]$, ainsi que l'inversion de
$\RightR{\lforall}$ qu'elle utilise, conservent les inférences
restantes dans le même ordre relatif et peuvent en supprimer.
Elles n'augmentent donc ni le nombre d'inférences quantifiées
ni le nombre d'inférences propositionnelles situées sous chacune.
L'inférence \RightR{\lforall} permutée avait une
\LeftR{\lif} au-dessous ; cette inférence propositionnelle
est maintenant au-dessus. En remplaçant $\pi_1$ dans $\pi$
par $\pi_1''$, l'ordre total diminue. Si l'inférence permutée
subsiste, son ordre a diminué d'au moins~$1$. Si elle disparaît,
sa contribution positive a été supprimée. La récurrence s'applique.

Le cas de \RightR{\lexists} est plus simple : aucune contraction
n'est nécessaire et cette règle n'impose pas de condition de
constante propre. Supposons, par exemple, que
$\RightR{\lexists}$ soit suivie de $\RightR{\land}$,
cette fois dans sa prémisse gauche. La sous-!!{proof}~$\pi_1$ est
\[
\AxiomC{}
\RightLabel{$\pi_2$}
\Deduce$\Gamma' \fCenter \Delta', \lexists[x][!A(x)], !A(t), !B$
\RightLabel{\RightR{\lexists}}
\UnaryInf$\Gamma' \fCenter \Delta', \lexists[x][!A(x)], !B$
\AxiomC{}
\RightLabel{$\pi_3$}
\Deduce$\Gamma' \fCenter \Delta', \lexists[x][!A(x)], !C$
\RightLabel{\RightR{\land}}
\BinaryInf$\Gamma' \fCenter \Delta', \lexists[x][!A(x)], !B \land !C$
\DisplayProof
\]
Remplaçons $\pi_1$ dans~$\pi$ par $\pi_1'$ :
\[
\AxiomC{}
\RightLabel{$\pi_2$}
\Deduce$\Gamma' \fCenter \Delta', \lexists[x][!A(x)], !A(t), !B$
\AxiomC{}
\RightLabel{$\pi_3'$}
\Deduce$\Gamma' \fCenter \Delta', \lexists[x][!A(x)], !A(t), !C$
\RightLabel{\RightR{\land}}
\BinaryInf$\Gamma' \fCenter \Delta', \lexists[x][!A(x)], !A(t), !B \land !C$
\RightLabel{\RightR{\lexists}}
\UnaryInf$\Gamma' \fCenter \Delta', \lexists[x][!A(x)], !B \land !C$
\DisplayProof
\]
La preuve~$\pi_3'$ s'obtient à partir de $\pi_3$ par affaiblissement
à droite avec~$!A(t)$. Après régularisation appropriée des constantes
propres, cet affaiblissement ajoute seulement $!A(t)$ à droite
dans les séquents de~$\pi_3$ ; il ne change donc pas l'ordre
des inférences quantifiées.

L'ordre des autres inférences quantifiées reste inchangé dans~$\pi'$.
Celui de l'inférence \RightR{\lexists} permutée avec
\RightR{\land} diminue de~$1$.
L'hypothèse de récurrence s'applique.
\end{proof}

\begin{prob}
Décrire les cas de la preuve du \olref{thm:midsequent} où
\LeftR{\lexists} aboutit à la prémisse de~\RightR{\lif},
et où \LeftR{\lforall} aboutit à la prémisse gauche
de~\LeftR{\lor}.
\end{prob}

\begin{proof}[Preuve du théorème de Herbrand]
Supposons que $\pi$ se termine par
$\Sequent \lexists[x][!A(x)]$, avec une matrice sans quantificateurs.
Par le \olref{thm:midsequent}, nous pouvons transformer $\pi$
en une !!{proof}~$\pi'$ possédant un séquent médian~$S$,
nécessairement de la forme
\[\Sequent
\lexists[x][!A(x)], !A(t_1), \dots, !A(t_n).\]
Comme toutes les inférences au-dessus de $S$ sont propositionnelles,
$\lexists[x][!A(x)]$ ne peut y être principale dans une inférence.
Elle n'est pas atomique et ne peut donc être la formule
d'identité d'un axiome. Puisque la matrice $!A(x)$ est sans
quantificateurs, la formule existentielle ne peut pas non plus
être une sous-formule propre de $!A(t_1)$, ni d'une autre instance,
et elle ne peut être active au-dessus de~$S$.

La suppression des occurrences de $\lexists[x][!A(x)]$ dans la
sous-!!{proof} se terminant en~$S$ donne donc une preuve correcte
de $\Sequent !A(t_1), \dots, !A(t_n)$.
Par $n-1$ inférences $\RightR{\lor}$, nous obtenons une preuve de
$\Sequent !A(t_1) \lor \dots \lor !A(t_n)$.
\end{proof}

% \begin{prob}
% Trouver une disjonction de Herbrand de $\lexists[x][\lforall[y][(!A(x) \lif
% !A(y))]]$ en construisant une preuve sans coupure de $\Sequent
% \lexists[x][\lforall[y][(!A(x) \lif !A(y))]]$, puis en appliquant
% la transformation de \olref{thm:midsequent} pour obtenir un
% séquent médian si nécessaire.
% \end{prob}














\end{document}
